\section{Experimental Setup}

\subsection{Model, data, and training horizon}

The main workload is an 8B-parameter Llama-family decoder with RMSNorm,
grouped-query attention, rotary position embeddings (RoPE), and a gated
feed-forward network \cite{dubey2024llama3}. The sequence length is 8,192 and
the global batch is 512 sequences, or 4,194,304 tokens per update. A
160B-token run ends at step 38,147. Unless stated otherwise, the named
low-precision route applies to every body linear in all 32 transformer blocks.

All main custom routes use a BF16 output projection followed by
Inductor-compiled regular cross entropy. The optimizer is AdamW with a
2,000-step linear warmup to a peak learning rate of $3\times10^{-4}$; the rate
remains constant after warmup over the reported horizon. Training losses are
compared at matched optimizer steps. Each long-run recipe is represented by
one trajectory, so the comparisons do not estimate seed-to-seed variation.
Table~\ref{tab:base-training-recipe} summarizes the shared configuration; the
appendix records execution-specific details.

\begin{table}[htbp]
  \centering
  \footnotesize
  \caption{Shared Llama-8B training configuration.}
  \label{tab:base-training-recipe}
  \begin{tabularx}{\textwidth}{@{}lY@{}}
    \toprule
    Field & Effective contract \\
    \midrule
    Model & 32 blocks; width 4,096; FFN 14,336; 32 Q heads; 8 KV heads; head dimension 128; vocabulary 128,256 \\
    Position / norm & Llama-3.1 RoPE with base $500{,}000$, $8\times$ scaling, low/high frequency factors 1 and 4, and original length 8,192; RMSNorm $\epsilon=10^{-5}$ \\
    Data / tokenizer & Shuffled OLMo Mix 1124 \cite{olmo2025olmo2}; Meta-Llama-3.1-8B tokenizer \\
    Optimizer & Fused AdamW, $(\beta_1,\beta_2)=(0.9,0.95)$, $\epsilon=10^{-8}$, weight decay 0.1, global gradient-norm clip 1.0 \\
    LR schedule & 2,000-step linear warmup to $3\times10^{-4}$; constant thereafter over the reported horizon \\
    Batch / horizon & Global batch 512; 4,194,304 tokens/update; 38,147 updates; 160,000,114,688 scheduled tokens \\
    State / output & BF16 parameters and non-FP4 operations; FP32 master/optimizer state; BF16 output projection; execution-specific regular CE \\
    \bottomrule
  \end{tabularx}
\end{table}

\subsection{Recipe ledger}

Tables~\ref{tab:recipe-ledger} and~\ref{tab:recipe-state-ledger} state the
numerical recipe for every reported route. ``Row SR'' means stochastic
rounding only for the row-oriented gradient consumed by Dgrad. Activations,
weights, scales, and the Wgrad-oriented gradient remain deterministic in those
custom routes. A Wgrad Hadamard transform always acts on $X$ and $dY$ as a
pair; it never transforms the stored weight.

\begin{table}[htbp]
  \centering
  \footnotesize
  \caption{Per-contraction numerical recipe. 2D16 and 2D32 denote one
  orientation-consistent $16\times16$ or $32\times32$ weight encoding.
  NV denotes global NVFP4, LC denotes CTA-local NVFP4, and MX denotes MXFP4.
  The 27/5 hybrid changes format by depth; the operand hybrid changes it by
  contraction inside every block.}
  \label{tab:recipe-ledger}
  \begin{tabularx}{\textwidth}{@{}YYYYYY@{}}
    \toprule
    Route & Body assignment & Fprop & Dgrad & Wgrad & Weight representation(s) \\
    \midrule
    BF16 & blocks 1--32 & BF16 & BF16 & BF16 & BF16 master \\
    TE-native \nvfp{} in all blocks & blocks 1--32 & global NV & global NV & global NV + TE-managed H16 & TE 2D16 \\
    TE \nvfp{} with four final BF16 blocks & NV blocks 1--28; BF16 29--32 & NV 1--28; BF16 29--32 & NV 1--28; BF16 29--32 & NV+TE-managed H16 1--28; BF16 29--32 & TE 2D16 / BF16 \\
    Global \nvfp{} v5 & blocks 1--32 & global NV & global NV & global NV & custom native 2D16 \\
    \mxfp{} + row-SR & blocks 1--32 & MX & MX & MX & custom 2D32 \\
    \localcta{}-v4 & blocks 1--32 & localCTA & localCTA & localCTA & custom 2D16 \\
    \localcta{}+fixed H16 & blocks 1--32 & localCTA & localCTA & localCTA + paired fixed-sign H16 & custom 2D16 \\
    27/5 depth hybrid & LC blocks 1--27; MX 28--32 & LC 1--27; MX 28--32 & LC 1--27; MX 28--32 & LC 1--27; MX 28--32 & LC 2D16; MX 2D32 \\
    \mxfp{}+row-SR+fixed H32 & blocks 1--32 & MX & MX & MX + paired fixed-sign H32 & custom 2D32 \\
    Operand hybrid (plain H16) & every block & MX & localCTA & MX + paired plain H16 & MX 2D32 + LC 2D16 encodings \\
    Operand hybrid (fixed H32) & every block & MX & localCTA & MX + paired fixed-sign H32 & MX 2D32 + LC 2D16 encodings \\
    \bottomrule
  \end{tabularx}
\end{table}

\begin{table}[htbp]
  \centering
  \footnotesize
  \caption{Rounding and precision policy. Every row uses a BF16 output
  projection. The machine-readable ledger records the full stochastic
  namespaces and the recovered seed metadata.}
  \label{tab:recipe-state-ledger}
  \begin{tabularx}{\textwidth}{@{}lYYY@{}}
    \toprule
    Route & Gradient rounding & Hadamard contract & BF16 body / head \\
    \midrule
    BF16 & n/a & none & all / BF16 \\
    TE-native \nvfp{} in all blocks & TE gradient SR & TE-managed H16 on Wgrad & none / BF16 \\
    TE \nvfp{} + four final BF16 blocks & TE gradient SR & TE-managed H16 on NV Wgrad & last four / BF16 \\
    Global \nvfp{} v5 & row data SR & none & none / BF16 \\
    \mxfp{} + row-SR & row data SR & none & none / BF16 \\
    \localcta{}-v4 & row data SR & none & none / BF16 \\
    \localcta{}+fixed H16 & row data SR & paired fixed-sign H16 on Wgrad $X,dY$ & none / BF16 \\
    27/5 depth hybrid & row data SR in both & none & none / BF16 \\
    \mxfp{}+row-SR+fixed H32 & row data SR & paired fixed-sign H32 on Wgrad $X,dY$ & none / BF16 \\
    Operand hybrid (plain H16) & LC row data SR & plain H16 ($D=I$) on Wgrad & none / BF16 \\
    Operand hybrid (fixed H32) & LC row data SR & paired fixed-sign H32 on Wgrad $X,dY$ & none / BF16 \\
    \bottomrule
  \end{tabularx}
\end{table}
\FloatBarrier

``TE-managed'' means that Transformer Engine controls the H16 transform. We do
not assign that route the stronger fixed-sign contract used by the custom
kernels.

All custom training routes retain BF16 parameters and FP32 optimizer state. The final column
in Table~\ref{tab:recipe-ledger} describes the transient FP4 representation
consumed by Tensor Cores. The operand hybrid emits two encodings from one BF16
weight: MX 2D32 for Fprop and CTA-local 2D16 for Dgrad. Wgrad does not consume
the stored weight; it quantizes $X$ and $dY$ with MXFP4. This is a
deliberate contraction-specific recipe rather than two layouts of one shared
FP4 weight.

Appendix~\ref{app:reproducibility} gives the full experiment-disposition
ledger and the checkpoint criteria for each result panel.

\subsection{End-to-end performance measurements}

The performance probes use the same accelerator type, sequence length, global
batch, seed, BF16 output projection, compiled CE, warmup rule, metric window,
and distributed telemetry. Each route uses its selected local batch,
activation-checkpoint policy, and distributed schedule, so the panel compares
complete end-to-end configurations rather than one isolated implementation
change. Tables report cross-device aggregates. Local batch (LB) and gradient
accumulation (GA) remain explicit because they affect occupancy, memory
lifetime, and the exposed output-layer cost.

We use the dense-model training FLOP convention implemented by TorchTitan. For
total parameter count $N$, embedding parameter count
$N_{\mathrm{emb}}$, number of layers $L$, query-head count $H$, head dimension
$d_h$, and sequence length $S$, it gives
\begin{equation}
  F_{\mathrm{model/token}}
  =6\bigl(N-N_{\mathrm{emb}}\bigr)+12LHd_hS.
  \label{eq:model-flops}
\end{equation}
The first term counts the forward and backward dense parameter operations;
the second counts the attention contractions without charging recomputation.
Here $N=8{,}030{,}261{,}248$, $N_{\mathrm{emb}}=525{,}336{,}576$, $L=32$,
$H=32$, $d_h=128$, and $S=8{,}192$, giving
$F_{\mathrm{model/token}}=57{,}914{,}449{,}920$ FLOPs/token.

For a route with per-device token rate $R$, we report BF16 model FLOP
utilization (BF16 MFU):
\begin{equation}
  \mathrm{MFU}_{\mathrm{BF16}}=100\frac{F_{\mathrm{model/token}}R}
  {F_{\mathrm{peak,BF16}}},
  \label{eq:mfu}
\end{equation}
using Equation~\ref{eq:model-flops} and the hardware's maximum dense BF16
Tensor Core throughput. NVIDIA specifies 360 PFLOP/s of sparse FP16/BF16
throughput for a 72-GPU GB200 NVL72 and states that dense throughput is half
the sparse value. Dividing by 72 GPUs and by two gives the maximum dense BF16
throughput used here, $F_{\mathrm{peak,BF16}}=2.5\times10^{15}$ FLOP/s/GPU,
or 2,500 BF16 TFLOP/s/GPU \cite{nvidia2025blackwell}. Tokens/s/GPU is the
directly observed metric. We recompute BF16 MFU from each unrounded throughput
aggregate with this peak. For an FP4 route, BF16 MFU is BF16-equivalent model
work divided by the BF16 peak; it measures effective training throughput, not
literal FP4 Tensor Core occupancy. It can therefore exceed 100\%.
Recomputing from the rounded token rate printed in a table can differ in the
last decimal place.

\subsection{Training and checkpoint evaluation}

Training curves use exact logged optimizer steps. Smoothed curves use an
exponential moving average only for visualization. Custom-route endpoints are
raw observations; Appendix~\ref{app:reproducibility} defines the BF16 and
TE-native reference curves.

Validation uses a fixed external stream constructed from an 82/18 DataComp-LM
(DCLM) and OLMo-without-DCLM source mixture
\cite{li2024datacomplm,olmo2025olmo2}. It contains 768
sequences and 6,291,456 next-token targets. Documents are deduplicated, mixed,
and shuffled before packing so a device rank cannot represent only one source
or cursor range. Every checkpoint consumes the same tokenized sequences. We
call this an external validation stream rather than a held-out split because
zero overlap with every historical training trajectory cannot be proven.
Before route comparison, converted and native evaluation must agree at the
full 8,192-token context, and a late BF16 checkpoint must improve over an early
one by more than five sequence-level standard errors.

Every downstream comparison uses the literal step-38,000 checkpoint and BF16
inference. Converted models must match TorchTitan's RoPE and RMSNorm semantics
before scoring. The evaluation panel follows the
NVIDIA-aligned subset of MMLU, HellaSwag, WinoGrande, and ARC-Challenge
\cite{hendrycks2021mmlu,gu2024olmes,nvidia2025nvfp4}. Shot counts and reported metrics appear with the
results.
